\[ %%%%%%%%%%%%%%%%%%%%%%%%%%%%%%%%%%%%%%%%%%%%%%%%%%%%%%%%%%%%%%%%%%%%%%%%%%%%%%%% \section{Mathematical Formulation and Learning Objective} \label{sec:mathematical_formulation} %%%%%%%%%%%%%%%%%%%%%%%%%%%%%%%%%%%%%%%%%%%%%%%%%%%%%%%%%%%%%%%%%%%%%%%%%%%%%%%% This section provides a formal mathematical formulation of the PHYSGEN framework. We define the problem of long-horizon prediction for nonlinear dynamical systems, introduce the physics-guided graph evolution operator, and formulate the learning objective that combines predictive accuracy, physical consistency, and stability regularization. %%%%%%%%%%%%%%%%%%%%%%%%%%%%%%%%%%%%%%%%%%%%%%%%%%%%%%%%%%%%%%%%%%%%%%%%%%%%%%%% \subsection{Problem Definition} \label{sec:problem_definition} %%%%%%%%%%%%%%%%%%%%%%%%%%%%%%%%%%%%%%%%%%%%%%%%%%%%%%%%%%%%%%%%%%%%%%%%%%%%%%%% We consider a nonlinear dynamical system whose state evolves over continuous or discrete time according to an unknown physical evolution process. The general formulation is: \begin{equation} X_{t+1} = \mathcal{F}(X_t,\Theta), \end{equation} where: \begin{itemize} \item $X_t$ represents the system state at time $t$; \item $\mathcal{F}$ denotes the unknown physical evolution operator; \item $\Theta$ represents system-specific physical parameters. \end{itemize} The objective of PHYSGEN is to learn an approximation of the evolution operator: \begin{equation} \hat{\mathcal{F}}_\theta \approx \mathcal{F}, \end{equation} such that future system states can be predicted: \begin{equation} \hat{X}_{t+k} = \hat{\mathcal{F}}_\theta^k(X_t). \end{equation} The central challenge considered in this work is long-horizon prediction, where the forecasting horizon $k$ is significantly larger than the one-step prediction interval. %%%%%%%%%%%%%%%%%%%%%%%%%%%%%%%%%%%%%%%%%%%%%%%%%%%%%%%%%%%%%%%%%%%%%%%%%%%%%%%% \subsubsection{Learning from Physical Trajectories} %%%%%%%%%%%%%%%%%%%%%%%%%%%%%%%%%%%%%%%%%%%%%%%%%%%%%%%%%%%%%%%%%%%%%%%%%%%%%%%% Given a dataset of observed system trajectories: \begin{equation} \mathcal{D} = \{ X_0,X_1,\dots,X_T \}, \end{equation} the model receives a sequence of physical states and learns the underlying dynamical transition: \begin{equation} (X_t,X_{t+1},...,X_{t+\tau}) \rightarrow \hat{X}_{t+\tau+1}. \end{equation} Unlike purely data-driven forecasting methods, PHYSGEN assumes that the observed trajectory is generated under physical constraints: \begin{equation} X_t\in\mathcal{M}_{phys}, \end{equation} where $\mathcal{M}_{phys}$ denotes the manifold of physically admissible states. Therefore, the learning objective is not only to minimize prediction error but also to preserve the structure of the physical state space. %%%%%%%%%%%%%%%%%%%%%%%%%%%%%%%%%%%%%%%%%%%%%%%%%%%%%%%%%%%%%%%%%%%%%%%%%%%%%%%% \subsubsection{Graph-Based System Representation} %%%%%%%%%%%%%%%%%%%%%%%%%%%%%%%%%%%%%%%%%%%%%%%%%%%%%%%%%%%%%%%%%%%%%%%%%%%%%%%% The physical state is represented as a dynamic graph: \begin{equation} G_t=(V_t,E_t,X_t,P_t), \end{equation} where: \begin{itemize} \item $V_t$ represents physical entities; \item $E_t$ represents interactions; \item $X_t$ contains node states; \item $P_t$ contains physical parameters. \end{itemize} The graph representation defines the mapping: \begin{equation} \mathcal{G}:X_t\rightarrow G_t. \end{equation} The objective is therefore reformulated as learning a graph dynamical operator: \begin{equation} G_{t+1} = \Phi_\theta(G_t,P_t). \end{equation} %%%%%%%%%%%%%%%%%%%%%%%%%%%%%%%%%%%%%%%%%%%%%%%%%%%%%%%%%%%%%%%%%%%%%%%%%%%%%%%% \subsubsection{Long-Horizon Forecasting Objective} %%%%%%%%%%%%%%%%%%%%%%%%%%%%%%%%%%%%%%%%%%%%%%%%%%%%%%%%%%%%%%%%%%%%%%%%%%%%%%%% The main prediction objective is defined as: \begin{equation} \min_\theta \mathbb{E} [ \| X_{t+k} - \hat{X}_{t+k} \|^2 ], \end{equation} where $k$ represents the prediction horizon. However, minimizing only short-term prediction error may result in unstable trajectories during recursive rollout. Therefore, PHYSGEN introduces additional physical and stability objectives: \begin{equation} \theta^* = \arg\min_\theta \mathcal{L}_{PHYSGEN}. \end{equation} The complete learning objective is defined as: \begin{equation} \boxed{ \mathcal{L}_{PHYSGEN} = \mathcal{L}_{data} + \lambda_p \mathcal{L}_{physics} + \lambda_s \mathcal{L}_{stability} } \end{equation} where: \begin{itemize} \item $\mathcal{L}_{data}$ measures prediction accuracy; \item $\mathcal{L}_{physics}$ enforces physical consistency; \item $\mathcal{L}_{stability}$ regulates long-term trajectory behavior. \end{itemize} This formulation defines PHYSGEN as a physics-guided learning framework in which prediction, physical reasoning, and dynamical stability are optimized jointly. \] 
